% !TEX root = ../main.tex
\chapter{Discovery: one template, many channels}\label{ch:discovery}
Analyse works on one span of one channel. Discovery takes what you built there and applies it at
scale: a saved \term{template}{template} run over several channels and many hours, or a
\term{seed-search}{seed search} for more of a shape you already have. Every Discovery run is
paired with its \term{null}{null}, so each result can be read against what chance would find on
the same stretch of recording. This is the workspace where a claim is made; Analyse is where a
method is tuned.

\begin{plainterms}[Why every run carries a null]
Suppose a detector finds 40 events in a day of recording. Is that a lot? You cannot tell until you
know how many it finds in a recording that has the same overall wiggle but no real events in it.
The app makes such stand-ins, called \term{surrogate}{surrogates}, by scrambling the recording in a
way that keeps its general character, runs the same template on them, and reports ``null expects
\(m\)'': the average count on the scrambled versions. ``\(40\) found, null expects 2'' is a finding;
``40 found, null expects 38'' is not.
\end{plainterms}

The four pages share one \emph{session}: a \textbf{scope} (one recording, one or more channels and a
section of hours) and a list of \textbf{runs} over it. \apppage{discovery.runs} is home: it applies
templates and shows where each run fires and how it scores. \apppage{discovery.seed} sets up a seed
search as a draft run in the same list. \apppage{discovery.compare} sets two runs side by side, and
\apppage{discovery.stages} pushes one place where they disagree through both chains, stage by stage.
All four read the bridge (\pagestatus{live}); the few controls that are still simulated are named
in each section. The held-out recording is refused on every Discovery route.

\begin{watchout}[Discovery runs write to the database]
In the project mode the researcher uses (\cref{ch:settings}), a Discovery run writes real rows:
runs, detections and the paired surrogate runs. That is the point, but it cannot be undone except
by \ui{Discard run}, which marks a run superseded.
\end{watchout}

% ---------------------------------------------------------------------------------------------
\pagesection{discovery.runs}{Runs}{Apply templates and seeds across channels}{live}
The Runs page applies saved templates to the scope, runs them, and shows the results: where each
run fires, a scoreboard against the human record and the null, and a browser to step through the
detections one at a time.

\pagefield{When you'd come here}
\begin{itemize}
  \item You saved a detection template in Analyse and want to run it over many channels.
  \item You want to know how a run scores: how many it found, how many a human judged, the
        precision, and how many the null expects.
  \item You want to send a run's detections to Review to be judged.
\end{itemize}

\screenshot{pages/discovery/discovery.runs--default}{A session: the scope strip at the top, the
runs list on the left, and on the right where each run fires per channel, the scoreboard and the
detection browser.}

\pagefield{Doing things here}

\subsection*{Set the scope}
\begin{enumerate}
  \item In the \ui{Scope} card, pick the \ui{recording}. Changing it asks first, because it clears
        the channels and marks every run stale.
  \item Add channels with \ui{+ channel} (a tick list with \ui{Select all}), or remove one with the
        cross on its chip. The scope keeps at least one channel. Strips show three channels at a
        time; the arrows beside them page through the rest.
  \item Set the section with the \ui{section} chip (from and to, in hours), or by dragging on the
        strips. The chip reads, for example, ``0--174 h · 174 h \(\times\) 3 ch''.
\end{enumerate}

\subsection*{Apply a template}
\begin{enumerate}
  \item Press \ui{Apply template} in the runs list. A window lists the saved templates that fit the
        scope, each with its chain drawn as a strip of glyphs, filterable and sortable.
  \item Tick one or more. The detail pane shows the template's locked parameters, its estimated
        cost per channel, and compute, disk and null tiles. A cost the bridge has not measured yet
        is left empty rather than guessed. For a template that starts with a filter, a band list
        (from Settings \textrightarrow{} Analysis defaults) can run it once per named band.
  \item The cost decides the button. Within the local limit (20 minutes by default, set in
        Settings \textrightarrow{} Compute \& HPC) the button is \ui{Add and run}; above it,
        \ui{Create SLURM script}.
\end{enumerate}
The toolbar's cost chip totals the pending runs (``\ldots{} local'' or ``\ldots{} routes to
cluster''); the toolbar button then reads \ui{Run \(n\)} or \ui{Create SLURM script for \(n\)}. The
null chip beside it says which surrogate method and how many draws the next template run and the
next seed search will carry; when no null method is available it says ``none'' and why.

\subsection*{Read where each run fires}
One row per run, grouped by channel, all on one time axis. In \ui{density} mode each 3-hour bin is
shaded by how many detections fell in it; in \ui{spans} mode each detection is one mark. The human
row shows reviewed hours as a grey underlay: recall can only be measured there. Click a bin to open
that run's first detection at or after it in the browser below. \ui{Open the map} goes to the
detection map across every dataset.

\subsection*{Read the scoreboard}
One row per run, expandable to its channels:
\begin{description}
  \item[found] detections the run made on the scope.
  \item[already judged] detections that had a human verdict before the run started.
  \item[reviewed, interesting, precision] how many fall where a human has judged, how many of those
        were judged interesting, and the share. Precision is shown as two figures, each with its
        rule on hover: \emph{containment} (over the windows a reviewer judged) and \emph{extent}
        (over the event-shaped human rows, matched by overlap).
  \item[recall · reviewed overlap] the share of human-marked events the run found, over the hours
        a human reviewed; the row states those hours, or says there are too few.
  \item[null expects / draws] the mean count on the run's surrogate draws, and how many draws that
        run made.
  \item[\(\times\) null] found divided by null expects. Empty when the null found nothing: there
        is no ratio to take, which is the best result there is.
\end{description}
After judging in Review, press \ui{Refresh after reviewing}: the scores are re-read from the
verdict tables, not recomputed.

\begin{watchout}[Precision alone rewards timidity]
A run that fires twice and is right both times has precision 1. Read precision with recall, the
found count and \(\times\) null, never alone.
\end{watchout}

\subsection*{Browse, review or discard a run}
\begin{enumerate}
  \item In \ui{Browse detections}, pick a run and a channel and step through its detections with
        the pager. Each shows its time, depth, score, any prior verdict, and which other runs found
        it too. Browsing writes nothing: verdicts are given in Review.
  \item \ui{Send \(n\) unjudged to Review} builds a Review queue over that run's detections that have
        no verdict yet (never its surrogates' detections) and offers \ui{Open Review}. Pressing it
        again reuses the same queue.
  \item \ui{Discard run} (after a confirmation) marks the run superseded. It writes no verdicts:
        its detections are \emph{not} marked uninteresting, because thousands of automatic ``no''
        verdicts would distort every later comparison between machine and human.
  \item A seed run offers a button to add its accepted matches to the Library entry it searched
        for.
  \item \ui{History} lists past Discovery runs; \ui{Open} brings one into this session.
\end{enumerate}

\pagefield{Where it leads}
\apppage{discovery.seed} (\ui{Seed search} in the runs list); \apppage{discovery.compare} (tick two
runs, A and B, then \ui{Compare}); \apppage{review.inspector} for the queue; Jobs for runs on the
cluster. \ui{Analyse events} is not wired yet.

\begin{watchout}[Simulated parts of this page]
The scope card's \ui{Preview on a 4 h sample} button plays a timed animation with example counts;
it measures nothing. The measured preview is the one in the \ui{Apply template} window.
\end{watchout}

\begin{underthehood}
Client: \code{webui/client/src/discovery/RunsPage.tsx}, \code{chrome.tsx} (\code{ScopeCard},
\code{RunsCard}, \code{NullChip}, \code{CostChip}, \code{SlurmModal}), \code{AddTemplateModal.tsx},
\code{session.ts}; reads in \code{webui/client/src/api/discovery.ts}. Bridge:
\code{webui/server/discovery.py} (\code{/api/discovery/session}, \code{/runs}, \code{/fires},
\code{/scoreboard}, \code{/detections}, \code{/templates/apply}, \code{/plan}, \code{/preview},
\code{/slurm}, \code{/runs/{run_key}/discard}, \code{/runs/{run_key}/review}). The fan-out and its
ceiling: \code{Working/discovery/fanout.py} (\code{ceiling_s}, Settings key
\code{limit.Discovery}).
\end{underthehood}

% ---------------------------------------------------------------------------------------------
\pagesection{discovery.seed}{Seed search}{Set up a seed run}{live}
A seed search finds more of a shape you already have. The \term{seed}{seed} is one example: a
Library entry, a span you judged in Review, or a span you took in Explore. The search slides it
along every channel in the scope and scores how alike each position is; the ones under a distance
cut are the matches. The page sets the search up as a draft run in the session's runs list.

\begin{plainterms}[How a seed search decides]
Picture laying a tracing of one shape over a long strip of recording and sliding it along, noting
at every position how far the strip is from the tracing. Most positions are far off; a few are
close. The \emph{distance} is that closeness score (0 is identical), and the \emph{cut} is the
largest distance you will still call a match. The null shows how close scrambled recordings come:
the sensible cut is where the null starts to produce matches too.
\end{plainterms}

\pagefield{When you'd come here}
\begin{itemize}
  \item You have one good example of an event and want to find others like it.
  \item Review or Explore sent you here with a seed already chosen (\ui{Seed search in Discovery}).
\end{itemize}

\screenshot{pages/discovery/discovery.seed--default}{The seed card with its provenance, the
parameters with the where-to-cut histogram (the null behind it), the distance profile on one
channel, and the apply bar.}

\pagefield{Doing things here}
\begin{enumerate}
  \item Check the \ui{Seed} card: the shape, and where it came from (recording, channel, samples,
        content hash). \ui{change seed} opens a picker filtered by kind (Library, Review, Explore
        selection) and paged through the Library.
  \item Choose what happens when the search is saved as a template: \emph{carry} keeps this exemplar
        with the template; \emph{rebind} asks for a new one each time it is applied.
  \item In \ui{Parameters}: the algorithm (MASS, a z-normalised distance); the window, locked at the
        seed's own length; an optional \ui{scale bank} that also searches stretched and squeezed
        copies of the seed; the \ui{exclusion zone} (matches closer than this to a better match are
        dropped); the \ui{match threshold}; and what to do \ui{on overlap}.
  \item Read the histogram of distances, with the null behind it, and drag the cut. A green tick
        marks the recommended cut, where the null starts to keep matches.
  \item The \ui{Distance profile} shows one channel and one zoomed view: the signal, the distance at
        every position with the cut, and the matches beneath. The \ui{Matches} card shows each match
        overlaid on the seed (both centred on their own mean, in millivolts), sorted by distance.
  \item Choose the run's null in the apply bar: \ui{from the preview} (the draws already made while
        you set up), \ui{rigorous · 200 paired runs}, or \ui{off}.
  \item Press \ui{Run seed search}. Or \ui{Save as template}, which stores the search with its cut,
        bank and exclusion zone (a taken name updates that template).
\end{enumerate}

\subsection*{The cost estimate, progress, and the cluster}
Before you press anything, the apply bar shows what the run will cost on this scope, for example
``about 4 min''. The estimate is \emph{work over a rate}: channels \(\times\) samples \(\times\) (one
search plus every null draw), divided by how much work this machine does per second. Until a search
has finished on this session the rate is an assumed one, and the estimate says so (``assumed
rate''); after that it is the rate measured on the last search.

While a search runs, a real progress bar shows the fraction done, the time so far, and about how
long is left; \ui{Cancel} stops it. When it finishes, \ui{Open in Runs} shows it in the runs list.

\begin{watchout}[Over ten minutes: run it on the cluster]
If the estimate is over the seed-search local limit (10 minutes unless Settings \textrightarrow{}
Compute \& HPC sets a Discovery limit), \ui{Run seed search} is disabled and \ui{Create SLURM
script} takes its place. The laptop is not meant for these.
\begin{enumerate}
  \item \ui{Create SLURM script} writes a \emph{spec} (the seed's own samples, the channels, the
        section, the parameters and the null) and an \texttt{sbatch} script that runs
        \texttt{python -m Working.discovery.seed\_job} on the cluster. The window shows the paths,
        the \texttt{sbatch} command, and where the result file is expected. The run row now reads
        \emph{on cluster}.
  \item Copy the spec and the script to the cluster and submit the job. It checkpoints after every
        channel and every ten null draws and resubmits itself while the result is incomplete, so a
        job stopped at the cluster's 20-minute wall time loses minutes, not the run.
  \item Bring the small result file back and press \ui{Import HPC result} (on the apply bar or in
        the window). The matches and the null draws enter the session's run.
\end{enumerate}
\end{watchout}

\pagefield{Where it leads}
\apppage{discovery.runs} to score the run and send its matches to Review; Library
\textrightarrow{} Templates for the saved search.

\begin{underthehood}
Client: \code{webui/client/src/discovery/SeedPage.tsx} (\code{SearchProgress},
\code{SeedSlurmModal}, \code{ImportResultButton}, \code{SaveTemplateModal}). Bridge:
\code{webui/server/discovery.py} \code{get_seed_estimate} (\code{_seed_estimate},
\code{_seed_work}: channels \(\times\) samples \(\times\) (1 + draws), divided by a rate measured
on the session or \code{DEFAULT_SEED_RATES}; ceiling \code{_seed_ceiling_s}, default
\code{SEED_CEILING_S} = 600\,s), \code{/api/discovery/seed/results}, \code{/seed/run},
\code{/seed/slurm}, \code{/seed/import}, \code{/seed/template}. The cluster job:
\code{Working/discovery/seed_job.py} (\code{build_spec}, \code{run_spec}, \code{write_job}), which
calls the same \code{seeded_search.candidates} and \code{null_distances} as the local preview. The
search block is \blockref{detection.seed_matches}.
\end{underthehood}

% ---------------------------------------------------------------------------------------------
\pagesection{discovery.compare}{Compare}{Template run A against seed run B}{live}
Compare sets two runs on the same scope side by side: what differs between their chains, where each
fires, how far their detections overlap, and every place they disagree. Either side can be the human
annotations instead of a run.

\pagefield{When you'd come here}
\begin{itemize}
  \item You ran two methods (a template and a seed search, or two settings of one template) and want
        to know where they agree.
  \item You want to measure a method against the human record.
  \item You want to judge only what the second method adds.
\end{itemize}

\screenshot{pages/discovery/discovery.compare--default}{Two runs compared: the chains aligned by
role, where A and B fire on one channel, the set overlap, and the disagreement stepper.}

\pagefield{Doing things here}
\begin{enumerate}
  \item On the Runs page tick one run as A and one as B and press \ui{Compare}. \ui{Swap A / B}
        exchanges them; \ui{all runs} goes back.
  \item \ui{What differs} lays the two chains out in columns by role, so stages doing the same job
        line up: identical stages grey, differing ones amber, an empty slot where a chain has no
        stage in that role. Click a stage for its parameters.
  \item \ui{Where A and B fire} draws one channel's signal clean, with A's detections, B's, and an
        agreement track beneath (grey both, blue only A, purple only B).
  \item \ui{Set overlap} counts \emph{only A}, \emph{both} and \emph{only B}, per channel. Two
        detections are the same event when their spans overlap by an intersection-over-union of at
        least 0.5 (Settings \textrightarrow{} Analysis defaults). Against the human record the rule
        is different, because most human labels are long review windows: a detection counts as
        \emph{both} when its human verdict is yes. A verdict split shows how much of each cell has
        been judged and accepted.
  \item \ui{Step through the disagreements} shows one window at a time: the signal, the span the
        firing run found, and the other run's own score there against its own threshold, so you see
        how close it came. Filter to only A or only B and page through.
  \item \ui{Send only-B unjudged to Review} queues what B found and A did not, for judging;
        \ui{Refresh after reviewing} re-reads the verdicts.
\end{enumerate}

\begin{plainterms}[Overlap by intersection over union]
Two spans of time overlap by some amount. Divide the shared part by the total stretch the two cover
together and you get a number from 0 (no overlap) to 1 (the same span). At 0.5 or more the app calls
them one event.
\end{plainterms}

\pagefield{Where it leads}
\apppage{discovery.stages} (\ui{Compare every stage}) to see why a disagreement happened; Review
for the queued remainder.

\begin{underthehood}
\code{webui/client/src/discovery/ComparePage.tsx}, \code{Divergence.tsx}, \code{Precision.tsx};
bridge \code{GET /api/discovery/compare}, \code{/compare/window}, \code{POST /compare/review} in
\code{webui/server/discovery.py}. Both runs are read from their stored results; nothing re-runs to
compare them.
\end{underthehood}

% ---------------------------------------------------------------------------------------------
\pagesection{discovery.stages}{Compare every stage}{One disagreement pushed through both runs}{live}
This page takes one window where A and B disagree and shows what every stage of each chain did
there, in two columns aligned by role, so you can see at which stage the two chains part ways.

\pagefield{When you'd come here}
From a disagreement on \apppage{discovery.compare}, when you want to know why one method fired and
the other did not.

\screenshot{pages/discovery/discovery.stages--default}{One window through both chains:
stage cards side by side, each badged identical, differs, A only or absent, and the table of the
numbers behind them.}

\pagefield{Doing things here}
\begin{enumerate}
  \item Step through disagreements with the pager. Each step runs both chains afresh on a window
        around the new disagreement; nothing is written. The page labels this ``40 s'', but the
        request asks the bridge for a 240-second window, and the ``about 1 s'' chip and the
        brief re-running veil are fixed, not measured.
  \item Read the stage cards. The step bar names the \emph{first stage whose output differs}. That
        is where the outputs diverge, not proof of what caused the disagreement.
  \item \ui{What each stage decided here} lists the numbers behind the pictures, one row per role;
        a tinted row is a role where the outputs differ. Click a stage for its parameters.
  \item \ui{Swap A / B}, or go back to Compare.
\end{enumerate}

\pagefield{Where it leads}
Back to \apppage{discovery.compare}. \ui{Open in Analyse} on a stage is not wired yet: it opens the
chain page without the run.

\begin{underthehood}
\code{webui/client/src/discovery/StagesPage.tsx}; bridge \code{GET /api/discovery/compare/stages}
in \code{webui/server/discovery.py}.
\end{underthehood}
